% One user's notation: what the shapes of tikz-tensors stand for in the figures
% of this repository. It is a convention, not part of the package -- the
% package draws shapes and lines and says nothing about what they mean; a
% project that means something else by a triangle keeps a file like this of
% its own. Every example loads it with % One user's notation: what the shapes of tikz-tensors stand for in the figures
% of this repository. It is a convention, not part of the package -- the
% package draws shapes and lines and says nothing about what they mean; a
% project that means something else by a triangle keeps a file like this of
% its own. Every example loads it with % One user's notation: what the shapes of tikz-tensors stand for in the figures
% of this repository. It is a convention, not part of the package -- the
% package draws shapes and lines and says nothing about what they mean; a
% project that means something else by a triangle keeps a file like this of
% its own. Every example loads it with % One user's notation: what the shapes of tikz-tensors stand for in the figures
% of this repository. It is a convention, not part of the package -- the
% package draws shapes and lines and says nothing about what they mean; a
% project that means something else by a triangle keeps a file like this of
% its own. Every example loads it with \input{notation}.
%
% Read top to bottom, it is the notation:
%   a function of position is a circle (or capsule) in the electron colour,
%   and its continuous argument a wavy leg in the same colour;
%   an array of numbers is a circle in the array colour -- a capsule, which is
%   a circle at one site's size and stretches to the sites it spans -- an
%   operator a plain rounded box -- no tensor is a square: a square
%   says nothing a circle does not, and turned on a lattice it reads as a
%   diamond;
%   left- and right-canonical tensors are triangles pointing the way the gauge
%   runs, each in its own colour, and the orthogonality center a green
%   diamond -- small when it sits on a bond; the gauge is said as the way
%   their bonds run (tn flux): into a canonical tensor from the far end of
%   the chain and out toward the center, and into the center from both
%   sides, so that an arrow on a bond (gauge) points toward the center;
%   an operator tensor on a chain is a neutral rounded box, a gate a light
%   yellow one, an environment block light purple;
%   a finite index is a plain line, and one whose direction is part of what is
%   being said -- which way the gauge runs -- carries an arrowhead;
%   two indices fused into one (a bond of dimension D times chi) are a double
%   line; a copy tensor, diagonal in all its indices, is an ink dot; a basis
%   vector closing an index is a small circle;
%   an isometry of a tree or a MERA is a triangle pointing up, toward the
%   root, in the left-canonical colour: off a chain there is no left and
%   right, only toward the center and away from it, and that is the colour
%   that already says `isometric toward the center'; both it and a MERA's
%   disentangler (a gate, flatter) are drawn flat, as a tree's tensors are;
%   the environment of a two-dimensional network, contracted from its
%   boundary, is the environment's colour: a corner as a rounded box, a side
%   as a circle;
%   in the tensor renormalization group, the tensor of the coarser lattice
%   is the array's shape a step paler, and the four halves a split leaves
%   are four tensors, so four colours: the two factors of one decomposition
%   two steps of one hue (warm for the halves above and below, green for
%   those to the left and right), the first factor the stronger -- one
%   colour, one tensor;
%   in a quantum circuit, a gate on one qubit is a gate, and a controlled NOT
%   is its control, an ink dot, joined to its target, a circled plus.
% The colours are the theme's (theme/tokens.toml). Shape carries the meaning
% on its own and colour repeats it, so a figure survives being printed in grey.
\tikzset{
  frame/.style      = {tn frame},
  tn gate/.append style = {fill=gate!25},
  tn env/.append style  = {fill=env!25},
  cont/.style       = {tn wavy, draw=ele},
  disc/.style       = {tn edge},
  gauge/.style      = {tn arrow},
  leg/.style        = {tn label},
  fused/.style      = {tn double},
  half1/.style      = {tn fill=warm2},
  half2/.style      = {tn fill=green2},
  half3/.style      = {tn fill=warm1},
  half4/.style      = {tn fill=green1}
}
% The tensors of the notation are types (\tntype): a tensor drawn in one says
% which, in \tnshow and in an error.
\tntype{fn}          {tn circle, draw=ele, fill=ele!12}
\tntype{fnwide}      {tn capsule, tn wide, draw=ele, fill=ele!12}
\tntype{fntall}      {tn capsule, tn tall, draw=ele, fill=ele!12}
\tntype{coef}        {tn capsule, tn fill=coef}
\tntype{coefwide}    {tn capsule, tn wide, tn fill=coef}
\tntype{coeftall}    {tn capsule, tn tall, tn fill=coef}
\tntype{op}          {tn rounded}
\tntype{opwide}      {op, tn wide}
\tntype{canl}        {tn triangle right, tn fill=canl, tn flux={left=in, right=out}}
\tntype{canr}        {tn triangle left, tn fill=canr, tn flux={right=in, left=out}}
\tntype{center}      {tn diamond, tn fill=center, tn flux={left=in, right=in}}
\tntype{centerbond}  {center, font=\scriptsize}
\tntype{mpo}         {tn rounded, tn fill=mpo}
\tntype{gate}        {tn gate}
\tntype{delta}       {tn dot}
\tntype{basis}       {tn circle, tn small}
\tntype{iso}         {tn triangle up, tn flat, tn fill=canl}
\tntype{disentangler}{gate, tn flat}
\tntype{isodown}     {tn triangle down, tn flat, tn fill=canl}
\tntype{corner}      {tn rounded, fill=env!25}
\tntype{side}        {tn circle, fill=env!25}
\tntype{coarse}      {tn capsule, tn fill=blue2}
\tntype{ctrl}        {tn dot}
\tntype{targ}        {tn oplus}

.
%
% Read top to bottom, it is the notation:
%   a function of position is a circle (or capsule) in the electron colour,
%   and its continuous argument a wavy leg in the same colour;
%   an array of numbers is a circle in the array colour -- a capsule, which is
%   a circle at one site's size and stretches to the sites it spans -- an
%   operator a plain rounded box -- no tensor is a square: a square
%   says nothing a circle does not, and turned on a lattice it reads as a
%   diamond;
%   left- and right-canonical tensors are triangles pointing the way the gauge
%   runs, each in its own colour, and the orthogonality center a green
%   diamond -- small when it sits on a bond; the gauge is said as the way
%   their bonds run (tn flux): into a canonical tensor from the far end of
%   the chain and out toward the center, and into the center from both
%   sides, so that an arrow on a bond (gauge) points toward the center;
%   an operator tensor on a chain is a neutral rounded box, a gate a light
%   yellow one, an environment block light purple;
%   a finite index is a plain line, and one whose direction is part of what is
%   being said -- which way the gauge runs -- carries an arrowhead;
%   two indices fused into one (a bond of dimension D times chi) are a double
%   line; a copy tensor, diagonal in all its indices, is an ink dot; a basis
%   vector closing an index is a small circle;
%   an isometry of a tree or a MERA is a triangle pointing up, toward the
%   root, in the left-canonical colour: off a chain there is no left and
%   right, only toward the center and away from it, and that is the colour
%   that already says `isometric toward the center'; both it and a MERA's
%   disentangler (a gate, flatter) are drawn flat, as a tree's tensors are;
%   the environment of a two-dimensional network, contracted from its
%   boundary, is the environment's colour: a corner as a rounded box, a side
%   as a circle;
%   in the tensor renormalization group, the tensor of the coarser lattice
%   is the array's shape a step paler, and the four halves a split leaves
%   are four tensors, so four colours: the two factors of one decomposition
%   two steps of one hue (warm for the halves above and below, green for
%   those to the left and right), the first factor the stronger -- one
%   colour, one tensor;
%   in a quantum circuit, a gate on one qubit is a gate, and a controlled NOT
%   is its control, an ink dot, joined to its target, a circled plus.
% The colours are the theme's (theme/tokens.toml). Shape carries the meaning
% on its own and colour repeats it, so a figure survives being printed in grey.
\tikzset{
  frame/.style      = {tn frame},
  tn gate/.append style = {fill=gate!25},
  tn env/.append style  = {fill=env!25},
  cont/.style       = {tn wavy, draw=ele},
  disc/.style       = {tn edge},
  gauge/.style      = {tn arrow},
  leg/.style        = {tn label},
  fused/.style      = {tn double},
  half1/.style      = {tn fill=warm2},
  half2/.style      = {tn fill=green2},
  half3/.style      = {tn fill=warm1},
  half4/.style      = {tn fill=green1}
}
% The tensors of the notation are types (\tntype): a tensor drawn in one says
% which, in \tnshow and in an error.
\tntype{fn}          {tn circle, draw=ele, fill=ele!12}
\tntype{fnwide}      {tn capsule, tn wide, draw=ele, fill=ele!12}
\tntype{fntall}      {tn capsule, tn tall, draw=ele, fill=ele!12}
\tntype{coef}        {tn capsule, tn fill=coef}
\tntype{coefwide}    {tn capsule, tn wide, tn fill=coef}
\tntype{coeftall}    {tn capsule, tn tall, tn fill=coef}
\tntype{op}          {tn rounded}
\tntype{opwide}      {op, tn wide}
\tntype{canl}        {tn triangle right, tn fill=canl, tn flux={left=in, right=out}}
\tntype{canr}        {tn triangle left, tn fill=canr, tn flux={right=in, left=out}}
\tntype{center}      {tn diamond, tn fill=center, tn flux={left=in, right=in}}
\tntype{centerbond}  {center, font=\scriptsize}
\tntype{mpo}         {tn rounded, tn fill=mpo}
\tntype{gate}        {tn gate}
\tntype{delta}       {tn dot}
\tntype{basis}       {tn circle, tn small}
\tntype{iso}         {tn triangle up, tn flat, tn fill=canl}
\tntype{disentangler}{gate, tn flat}
\tntype{isodown}     {tn triangle down, tn flat, tn fill=canl}
\tntype{corner}      {tn rounded, fill=env!25}
\tntype{side}        {tn circle, fill=env!25}
\tntype{coarse}      {tn capsule, tn fill=blue2}
\tntype{ctrl}        {tn dot}
\tntype{targ}        {tn oplus}

.
%
% Read top to bottom, it is the notation:
%   a function of position is a circle (or capsule) in the electron colour,
%   and its continuous argument a wavy leg in the same colour;
%   an array of numbers is a circle in the array colour -- a capsule, which is
%   a circle at one site's size and stretches to the sites it spans -- an
%   operator a plain rounded box -- no tensor is a square: a square
%   says nothing a circle does not, and turned on a lattice it reads as a
%   diamond;
%   left- and right-canonical tensors are triangles pointing the way the gauge
%   runs, each in its own colour, and the orthogonality center a green
%   diamond -- small when it sits on a bond; the gauge is said as the way
%   their bonds run (tn flux): into a canonical tensor from the far end of
%   the chain and out toward the center, and into the center from both
%   sides, so that an arrow on a bond (gauge) points toward the center;
%   an operator tensor on a chain is a neutral rounded box, a gate a light
%   yellow one, an environment block light purple;
%   a finite index is a plain line, and one whose direction is part of what is
%   being said -- which way the gauge runs -- carries an arrowhead;
%   two indices fused into one (a bond of dimension D times chi) are a double
%   line; a copy tensor, diagonal in all its indices, is an ink dot; a basis
%   vector closing an index is a small circle;
%   an isometry of a tree or a MERA is a triangle pointing up, toward the
%   root, in the left-canonical colour: off a chain there is no left and
%   right, only toward the center and away from it, and that is the colour
%   that already says `isometric toward the center'; both it and a MERA's
%   disentangler (a gate, flatter) are drawn flat, as a tree's tensors are;
%   the environment of a two-dimensional network, contracted from its
%   boundary, is the environment's colour: a corner as a rounded box, a side
%   as a circle;
%   in the tensor renormalization group, the tensor of the coarser lattice
%   is the array's shape a step paler, and the four halves a split leaves
%   are four tensors, so four colours: the two factors of one decomposition
%   two steps of one hue (warm for the halves above and below, green for
%   those to the left and right), the first factor the stronger -- one
%   colour, one tensor;
%   in a quantum circuit, a gate on one qubit is a gate, and a controlled NOT
%   is its control, an ink dot, joined to its target, a circled plus.
% The colours are the theme's (theme/tokens.toml). Shape carries the meaning
% on its own and colour repeats it, so a figure survives being printed in grey.
\tikzset{
  frame/.style      = {tn frame},
  tn gate/.append style = {fill=gate!25},
  tn env/.append style  = {fill=env!25},
  cont/.style       = {tn wavy, draw=ele},
  disc/.style       = {tn edge},
  gauge/.style      = {tn arrow},
  leg/.style        = {tn label},
  fused/.style      = {tn double},
  half1/.style      = {tn fill=warm2},
  half2/.style      = {tn fill=green2},
  half3/.style      = {tn fill=warm1},
  half4/.style      = {tn fill=green1}
}
% The tensors of the notation are types (\tntype): a tensor drawn in one says
% which, in \tnshow and in an error.
\tntype{fn}          {tn circle, draw=ele, fill=ele!12}
\tntype{fnwide}      {tn capsule, tn wide, draw=ele, fill=ele!12}
\tntype{fntall}      {tn capsule, tn tall, draw=ele, fill=ele!12}
\tntype{coef}        {tn capsule, tn fill=coef}
\tntype{coefwide}    {tn capsule, tn wide, tn fill=coef}
\tntype{coeftall}    {tn capsule, tn tall, tn fill=coef}
\tntype{op}          {tn rounded}
\tntype{opwide}      {op, tn wide}
\tntype{canl}        {tn triangle right, tn fill=canl, tn flux={left=in, right=out}}
\tntype{canr}        {tn triangle left, tn fill=canr, tn flux={right=in, left=out}}
\tntype{center}      {tn diamond, tn fill=center, tn flux={left=in, right=in}}
\tntype{centerbond}  {center, font=\scriptsize}
\tntype{mpo}         {tn rounded, tn fill=mpo}
\tntype{gate}        {tn gate}
\tntype{delta}       {tn dot}
\tntype{basis}       {tn circle, tn small}
\tntype{iso}         {tn triangle up, tn flat, tn fill=canl}
\tntype{disentangler}{gate, tn flat}
\tntype{isodown}     {tn triangle down, tn flat, tn fill=canl}
\tntype{corner}      {tn rounded, fill=env!25}
\tntype{side}        {tn circle, fill=env!25}
\tntype{coarse}      {tn capsule, tn fill=blue2}
\tntype{ctrl}        {tn dot}
\tntype{targ}        {tn oplus}

.
%
% Read top to bottom, it is the notation:
%   a function of position is a circle (or capsule) in the electron colour,
%   and its continuous argument a wavy leg in the same colour;
%   an array of numbers is a circle in the array colour -- a capsule, which is
%   a circle at one site's size and stretches to the sites it spans -- an
%   operator a plain rounded box -- no tensor is a square: a square
%   says nothing a circle does not, and turned on a lattice it reads as a
%   diamond;
%   left- and right-canonical tensors are triangles pointing the way the gauge
%   runs, each in its own colour, and the orthogonality center a green
%   diamond -- small when it sits on a bond; the gauge is said as the way
%   their bonds run (tn flux): into a canonical tensor from the far end of
%   the chain and out toward the center, and into the center from both
%   sides, so that an arrow on a bond (gauge) points toward the center;
%   an operator tensor on a chain is a neutral rounded box, a gate a light
%   yellow one, an environment block light purple;
%   a finite index is a plain line, and one whose direction is part of what is
%   being said -- which way the gauge runs -- carries an arrowhead;
%   two indices fused into one (a bond of dimension D times chi) are a double
%   line; a copy tensor, diagonal in all its indices, is an ink dot; a basis
%   vector closing an index is a small circle;
%   an isometry of a tree or a MERA is a triangle pointing up, toward the
%   root, in the left-canonical colour: off a chain there is no left and
%   right, only toward the center and away from it, and that is the colour
%   that already says `isometric toward the center'; both it and a MERA's
%   disentangler (a gate, flatter) are drawn flat, as a tree's tensors are;
%   the environment of a two-dimensional network, contracted from its
%   boundary, is the environment's colour: a corner as a rounded box, a side
%   as a circle;
%   in the tensor renormalization group, the tensor of the coarser lattice
%   is the array's shape a step paler, and the four halves a split leaves
%   are four tensors, so four colours: the two factors of one decomposition
%   two steps of one hue (warm for the halves above and below, green for
%   those to the left and right), the first factor the stronger -- one
%   colour, one tensor;
%   in a quantum circuit, a gate on one qubit is a gate, and a controlled NOT
%   is its control, an ink dot, joined to its target, a circled plus.
% The colours are the theme's (theme/tokens.toml). Shape carries the meaning
% on its own and colour repeats it, so a figure survives being printed in grey.
\tikzset{
  frame/.style      = {tn frame},
  tn gate/.append style = {fill=gate!25},
  tn env/.append style  = {fill=env!25},
  cont/.style       = {tn wavy, draw=ele},
  disc/.style       = {tn edge},
  gauge/.style      = {tn arrow},
  leg/.style        = {tn label},
  fused/.style      = {tn double},
  half1/.style      = {tn fill=warm2},
  half2/.style      = {tn fill=green2},
  half3/.style      = {tn fill=warm1},
  half4/.style      = {tn fill=green1}
}
% The tensors of the notation are types (\tntype): a tensor drawn in one says
% which, in \tnshow and in an error.
\tntype{fn}          {tn circle, draw=ele, fill=ele!12}
\tntype{fnwide}      {tn capsule, tn wide, draw=ele, fill=ele!12}
\tntype{fntall}      {tn capsule, tn tall, draw=ele, fill=ele!12}
\tntype{coef}        {tn capsule, tn fill=coef}
\tntype{coefwide}    {tn capsule, tn wide, tn fill=coef}
\tntype{coeftall}    {tn capsule, tn tall, tn fill=coef}
\tntype{op}          {tn rounded}
\tntype{opwide}      {op, tn wide}
\tntype{canl}        {tn triangle right, tn fill=canl, tn flux={left=in, right=out}}
\tntype{canr}        {tn triangle left, tn fill=canr, tn flux={right=in, left=out}}
\tntype{center}      {tn diamond, tn fill=center, tn flux={left=in, right=in}}
\tntype{centerbond}  {center, font=\scriptsize}
\tntype{mpo}         {tn rounded, tn fill=mpo}
\tntype{gate}        {tn gate}
\tntype{delta}       {tn dot}
\tntype{basis}       {tn circle, tn small}
\tntype{iso}         {tn triangle up, tn flat, tn fill=canl}
\tntype{disentangler}{gate, tn flat}
\tntype{isodown}     {tn triangle down, tn flat, tn fill=canl}
\tntype{corner}      {tn rounded, fill=env!25}
\tntype{side}        {tn circle, fill=env!25}
\tntype{coarse}      {tn capsule, tn fill=blue2}
\tntype{ctrl}        {tn dot}
\tntype{targ}        {tn oplus}

